\documentclass[a4paper]{report}

%====================== PACKAGES ======================

\usepackage[french]{babel}
\usepackage[utf8x]{inputenc}
%pour gérer les positionnement d'images
\usepackage{float}
\usepackage{amsmath}
\usepackage{graphicx}
\usepackage[colorinlistoftodos]{todonotes}
\usepackage{url}
%pour les informations sur un document compilé en PDF et les liens externes / internes
\usepackage{hyperref}
%pour la mise en page des tableaux
\usepackage{array}
\usepackage{tabularx}
%pour utiliser \floatbarrier
%\usepackage{placeins}
%\usepackage{floatrow}
%espacement entre les lignes
\usepackage{setspace}
%modifier la mise en page de l'abstract
\usepackage{abstract}
%police et mise en page (marges) du document
\usepackage[T1]{fontenc}
\usepackage[top=2cm, bottom=2cm, left=2cm, right=2cm]{geometry}
%Pour les galerie d'images
\usepackage{subfig}

%====================== INFORMATION ET REGLES ======================

%rajouter les numérotation pour les \paragraphe et \subparagraphe
\setcounter{secnumdepth}{4}
\setcounter{tocdepth}{4}

\hypersetup{                            % Information sur le document
pdfauthor = {Sow Mamadou,},          % Auteurs
pdftitle = {Rapport de stage -
            Sujet du Projet},           % Titre du document
pdfsubject = {Mémoire de Projet},       % Sujet
pdfkeywords = {Tag1, Tag2, Tag3, ...},  % Mots-clefs
pdfstartview={FitH}}                    % ajuste la page à la largueur de l'écran
%pdfcreator = {MikTeX},% Logiciel qui a crée le document
%pdfproducer = {}} % Société avec produit le logiciel

%======================== DEBUT DU DOCUMENT ========================

\begin{document}

%régler l'espacement entre les lignes
\newcommand{\HRule}{\rule{\linewidth}{0.5mm}}

%page de garde
\begin{titlepage}
\begin{center}

% Upper part of the page. The '~' is needed because only works if a paragraph has started.
\includegraphics[width=0.35\textwidth]{./logo_lyon1_2.jpg}~\\[0.2cm]

\textsc{\LARGE Université Claude Bernard Lyon 1}\\
Adresse: \textsc{43, Boulevard du 11 Novembre  \\ Villeurbanne 69100}\\[1cm]

\textsc{\Large }\\[0.5cm]

% Title
{\huge \bfseries Rapport De Stage\\}

\HRule \\[0.4cm]

{\huge \bfseries \Large Développement suivant un modèle MVC de La définition de permissions sur la console d'administration en fonction du profil des utilisateurs.} \\[0.4cm]
 %Definition des droits sur une console d'administration en fonction du profil des utilisateurs%

\HRule \\[1.5cm]

%logo entreprise
\includegraphics[width=0.35\textwidth]{./secuserve_logo.png}~\\[0.2cm]

%nom de l'entreprise
\textsc{\LARGE Entreprise Secuserve}\\
Adresse: \textsc{74, Bd du 11 Novembre 1918 \\ 69100 VILLEURBANNE}\\[2.5cm]

% Author and supervisor
\begin{minipage}{0.4\textwidth}
\begin{flushleft} \large
\emph{Stagiaire}\\
Nom: \textbf{\textsc{Sow}}\\
 Prénom: \textbf{\textsc{Mamadou}}\\
Numéro-étudiant: \textbf{\textsc{11414531}}\\
\textbf{\textsc{Licence 3 Informatique}}\\
Année-Scolaire: \textbf{\textsc{2016-2017}}
\end{flushleft}
\end{minipage}
\begin{minipage}{0.4\textwidth}
\begin{flushright} \large
 \textbf{\textsc{}}\\
\emph{Maitre de stage} \\
Nom: \textbf{\textsc{Fèvre  }}\\
Prénom: \textbf{\textsc{Pièrre   }}\\
\textbf{\textsc{Ingénieur DEVOPS}}\\
Lieu:   \textbf{\textsc{Lyon    }}\\
\end{flushright}
\end{minipage}




\vfill

% Bottom of the page
{\large \today}

\end{center}
\end{titlepage}

%page blanche
\newpage
~
%ne pas numéroter cette page
\thispagestyle{empty}
\newpage

%\renewcommand{\abstractname}{Remerciements}
\addcontentsline{toc}{chapter}{Remerciements}
\begin{abstract}
\hskip7mm
\begin{spacing}{1.3}
Tout d'abord je tiens à remercier l'ensemble du personnel de l'entreprise Secuserve pour votre accueil chaleureux. \\
Au présidant Mr Stéphane Bouché de m'avoir donner la chance d'exploiter et d'acquérir des nouvelles expériences au seins de son entreprise.\\
Je remercie également mon maitre de stage Mr Pièrre Fèvre de m'avoir choisi et d'avoir eu confiance en moi pour intégrer dans cette structure jeune, dynamique, motivé et passionné de l'informatique.\\
Mais aussi à Mr Lionel Fèvre le responsable R\&D que je respecte énormément, qui regorge plein de ressource et de motivation avec lequel je travaille et apprend plein de nouvelles choses dans ce projet.\\

Enfin je remercie mon tuteur Fabien Duchateau pour le temps qu'il accorde à cette expérience professionnel ma référente pédagogique Marie Lefevre qui à toujours été présente quand j'en ai eu besoin durant toute cette année scolaire 2016-2017 à l'université Lyon 1.   
\end{spacing}
\end{abstract}

\tableofcontents
\thispagestyle{empty}
\setcounter{page}{0}
%ne pas numéroter le sommaire

\newpage

%espacement entre les lignes d'un tableau
\renewcommand{\arraystretch}{1.5}

%====================== INCLUSION DES PARTIES ======================

~
\thispagestyle{empty}
%recommencer la numérotation des pages à "1"
\setcounter{page}{0}
\newpage

%

\begin{center}
%\vspace{5cm }
\section{ \colorbox{orange}{Introduction}}
\end{center}
\begin{spacing}{1.3}
Stage en Licence 3 Informatique d’une durée de 4 mois en développement web suivant le modèle MVC (PHP, JAVA-SCRIPT, BOOTSTRAP, JQUERY, XML etc...), dans l'entreprise Secuserve qui est éditeur des services de sécurité et de messagerie.\\
Mon stage à débuté le 02/05/2017 et se terminera le 31/08/2017 dans les locaux des Secuserve à Villeurbanne, Au-delà d'enrichir et de mettre en pratique mes connaissances acquises, j'ai pu confirmer mon souhait de devenir développeur web.\\ \\
L'objet du stage porte sur deux aspects en corrélation:
D'une part la définition de permissions sur la console d'administration en fonction du profil des utilisateurs issu d'une structure hiérarchique.\\
D'autre part la définition des règles de communication et de sécurité pour le filtrage des messages en sortie.\\
Les développements suivront le modèle mvc et pourront utiliser des outils tel-que: java script, jquery, bootstrap, etc..., tout développement devra être testé et documenté.\\
Travaillant avec le responsable de R\&D , j'apprends dans d'excellentes conditions en adoptant des méthodes de travail pratiques et efficaces.  \\

En vue de rendre compte de manière fidèle et analytique le temps passé au sein de cette entreprise depuis le mois de mai jusqu'à présent, je vais commencer par vous présenter l'environnement professionnel du stage en suite tous les aspects de mon travail jusqu'à maintenant et faire une mini présentation de la seconde partie de mon stage qui commence après le mois de juin en précisant les différentes missions et tâches que j'ai pu effectuer au sein du service. Je finirai par la conclusion  en résumant les apports de mon stage.
\end{spacing}

%\input{./existant.tex}
 
%\input{./besoins.tex}

%\input{./autre_partie.tex}

%\input{./resultats.tex}

%\input{./bilan.tex}

%\input{./annexes.tex}

\newpage

%récupérer les citation avec "/footnotemark"
\nocite{*}

%choix du style de la biblio
\bibliographystyle{plain}
%inclusion de la biblio
%\bibliography{bibliographie.bib}
%voir wiki pour plus d'information sur la syntaxe des entrées d'une bibliographie

\end{document}